\section{Positive expressions, identities, and finite zeros}
\label{u:sec:positive}

\begin{definition}[Positive fixed-limit expressions]\label{u:def:positive}
For a finite list of variables $x_1,\ldots,x_d$ valued in $\NN$, the input
language $\mathcal P_d$ is generated by natural constants, those variables,
addition, and $h_\lambda(A)=\psi(\lambda+A)$, for fixed closed nonzero limits
$\lambda<\varepsilon_0$. Only finitely many heads occur in any input.
Natural coefficient multiplication means repeated addition. There is no
subtraction inside a head, multiplication of two open values, variable
ordinal index, inverse control or additional orbit-height operation.
\end{definition}
A difference $A-B$ is allowed as a signed row to be compared; it is not a
new constructor available inside an input expression. Indices may be
arbitrarily large below $\varepsilon_0$, and nesting is unrestricted.

Three conclusions will be established. Universal equality is decidable in
$\mathcal P_d$ for any finite $d$. In $\mathcal P_1$, a nonidentity has only
effectively bounded finitely many zeros. Finally, in $\mathcal P_1$ its
eventual sign is effective. The third conclusion requires substantially
more than the second: a nonzero integer-valued sequence can change sign
infinitely often without ever being zero.


The greedy evaluator of Corollary~\ref{f:cor:closed-evaluation} gives,
for an infinite ordinal input of norm $t$, the exact decomposition
\begin{equation}
  \psi(\alpha)=f^L(t)+L,\qquad L\ge1,                 \label{u:display:1-3}
\end{equation}
where $L$ counts replacements until the first finite state.  In particular
all displayed closed constants are computable.  The control estimates and
controlled strictness used throughout are
\begin{align}
 f(v)&\ge2^v>v,& f(v+1)&\ge2^{f(v)+1}, \label{u:eq:gaps}\\
 \alpha<\beta\ \&\ \norm(\alpha)\le f(\norm(\beta))
 &\ \Longrightarrow\ \psi(\alpha)<\psi(\beta).
 \label{u:eq:strict}
\end{align}
The first line follows from Lemma~\ref{m:lem:growth}; the second is a
direct clause of the collapse definition.  Neither asserts unrestricted
ordinal monotonicity of the collapse.

\begin{lemma}[Sparse peaks]\label{u:lem:sparse}
For an infinite input of norm $t\ge2$, write $Z=f^L(t)=f(s)$ as in \eqref{u:display:1-3}.
Then $0<L\le s-1$, and
\begin{equation}
           Z<\psi(\alpha)=Z+L<2Z<f(s+1).             \label{u:display:1-6}
\end{equation}
Distinct peaks $f(s)$ and $f(u)$ with $u<s$ satisfy $f(u)\le f(s-1)$.
\end{lemma}
\begin{proof}
Each norm update increases the norm by at least one. Thus
$s=f^{L-1}(t)\ge t+L-1\ge L+1$. Since $f(s)>s$, the count is smaller than
$Z$. The final strict inequality follows from \eqref{u:eq:gaps}. The last
assertion uses strict increase of the integer map $f$.
\end{proof}

\subsection{A complete normal form for identities}
\begin{definition}[Collected recursive normal form]\label{u:def:nf}
A normal form consists of a constant $c\in\N$, a vector
$b\in\N^d$, and a finite dictionary of tags $(\lambda,A)$ with positive
natural coefficients. Each $A$ is itself a normal form with nonempty
variable support. Its value is
\begin{equation}
       c+\sum_i b_i x_i+\sum_{(\lambda,A)} a_{\lambda,A}h_\lambda(A).
                                                               \label{u:display:2-1}
\end{equation}
The dictionary graph is finite and acyclic. The height is zero if no tag
occurs, and otherwise one plus the largest argument height.
\end{definition}
Normalize bottom up. Addition merges dictionaries and coefficient vectors.
For a head, normalize its argument; if that argument is closed, evaluate
the entire call and return a constant. Otherwise keep the tag
$(\lambda,\NF(A))$. Indices are compared as actual canonical Cantor forms.
A fixed finite tail in a presented index is first moved into its argument;
a zero limit part is removed using $\psi(n)=n$.

Identical normal forms have equal values, by induction. Completeness is a
separation theorem and needs more than commutative algebra.

\begin{lemma}[Orbit cancellation at roots]\label{u:lem:roots}
Call a natural number an $f$-root if it is outside $\operatorname{range}(f)$.
If infinite inputs $\alpha,\beta$ have root norms $t,u$ and equal peaks,
then $\alpha=\beta$.
\end{lemma}
\begin{proof}
From $f^L(t)=f^J(u)$ cancel the smaller number of iterates, using injectivity.
If $L\ne J$, one of the two roots is a positive iterate of the other, a
contradiction. Hence $L=J$ and $t=u$. Formula \eqref{u:display:1-3} gives equal collapse
values. If the ordinals differed, order them; their common norm satisfies
$t\le f(t)$, so \eqref{u:eq:strict} would give distinct values.
\end{proof}
This proof is useful because it avoids any assertion that terminal-peak
maps are monotone in arbitrary ordinal inputs.

\begin{theorem}[Identity normal form]\label{u:thm:identity}
For $E,F\in\mathcal P_d$, the identity $E=F$ holds at every positive
valuation exactly when $\NF(E)=\NF(F)$. If the forms differ, one can
effectively specify a positive countervaluation of the form
\begin{equation}
           x_i=w_i(f(k)+1),\qquad w_i,k\in\NN.        \label{u:display:2-2}
\end{equation}
Consequently universal equality in $\mathcal P_d$ is decidable.
\end{theorem}
\begin{proof}
If both expressions are closed, evaluate them; when they differ take
$k=1$ and every $w_i=1$ for the required countervaluation, with the
empty tuple allowed when $d=0$.  Otherwise let $\mathcal S$ contain
the two root normal forms, zero, and all recursively
reachable tag arguments. Let $M$ be the largest raw-variable coefficient
in this finite set, $B=\max(2,M+1)$, and $w_i=B^{i-1}$.
Uniqueness of finite base-$B$ notation shows that different coefficient
vectors in $\mathcal S$ have different weighted slopes $b\cdot w$.
On the ray $x_i=w_i n$, any open positive expression is at least $n$:
this follows by induction from positivity and $h_\lambda(q)\ge q$.

We first make every reachable open atom's initial norm a root on the ray
$n=f(k)+1$. For a tag $(\lambda,A)$ write that norm as
\begin{equation}
 t= A_0+Dn+\sum_j a_j h_{\lambda_j}(A_j),\quad
 A_0=\norm(\lambda)+c_A\ge2,\quad D=b_A\cdot w,
 \quad C=\sum_j a_j.                                 \label{u:display:2-3}
\end{equation}
If $C=0$, openness forces $D\ge1$. For $k\ge\max(4,A_0+2D)$,
\[
          f(k)<t\le(A_0+2D)f(k)<f(k+1),
\]
by \eqref{u:eq:gaps}. If $C>0$, let $Y=f(s)$ be the largest child peak.
Every child is open, so $Y\ge f(n)\ge2^n$. For
$n\ge\max(4,A_0,D+1,2C+1)$ the affine part is at most $n^2\le Y$.
The largest child contributes more than $Y$, while every child is below
$2Y$ by Lemma~\ref{u:lem:sparse}. Therefore
\begin{equation}
                 Y<t<(2C+1)Y<f(s+1).                \label{u:display:2-4}
\end{equation}
This too makes $t$ a root. All these bounds are finite syntax data.

Induct on the height of forms in $\mathcal S$. Distinct affine forms are
unequal for all sufficiently large $n$, because their slopes differ or
their constants differ. Suppose distinct lower-height forms have already
been separated. Two distinct tags at the next height cannot have equal
peaks: Lemma~\ref{u:lem:roots} would identify the ordinals
$\lambda+A(x)$ and $\mu+B(x)$, hence their limit parts and finite tails.
The indices would agree and the different lower-height arguments would
have equal values, contradicting induction.

Now cancel identical tags in the difference of two distinct forms at this
height. If tags remain, write the row on the weighted ray as
\begin{equation}
                \sum_j c_j h_{\lambda_j}(A_j)+pn+q,  \label{u:display:2-5}
\end{equation}
where $c_j\ne0$ and the remaining tags are distinct. Let $Y=f(s)$ be the
largest peak. Its coefficient has absolute value at least one. All other
peaks are at most $U=f(s-1)$, and their full values are below $2U$.
Because the largest tag is open, $s\ge n$, while $U\ge s$. Set
$K=2\sum_j|c_j|+|p|+|q|$. The absolute value of all lower contributions
is at most $KU$. Once $k\ge\max(4,K)$, the leading peak satisfies
$Y\ge2^{U+1}>KU$. Thus the row is nonzero, whichever tag is largest.
No algorithm for determining that largest tag has been assumed.

For explicit effectivity take the maximum of the root bounds above and,
for every pair in $\mathcal S$, the displayed row bound $K$; for an affine
row $pn+q$ with $p\ne0$, also include $|q|+1$. This single computable
$k$ suffices at every induction height. Normalization and all finite maxima
terminate, proving the theorem.
\end{proof}

\begin{example}[The exact boundary of positivity]
The collected forms of $h_\omega(x+y)+x$ and $x+h_\omega(y+x)$ agree.
In contrast the two tags in $h_\omega(x)+h_\omega(y)$ remain distinct;
the theorem decides an identity, not whether they happen to coincide at
some valuation. Inner subtraction cannot be added silently. The literal
finite-$\omega$ evaluation gives the exact identity
\[
 h_{2\omega}(n)=h_\omega(h_\omega(n+2)-3)+1
              =f(f(n+4))+2.
\]
The right side's subtraction is forbidden by Definition~\ref{u:def:positive}.
The free normal form is not asserted for that larger language.
\end{example}

\subsection{An effective finite-zero bound in one variable}
From now on every open input depends on the same $x\in\NN$.
This hypothesis permits a stronger separation statement along the entire
tail, without selecting the sparse ray \eqref{u:display:2-2}.

\begin{lemma}[Common-band or root separation]\label{u:lem:oneinput}
For any finite collection of normal forms in $\mathcal P_1$, a computable
$N$ has the following properties for every $x\ge N$.
Every affine open atom input lies in $[x,f(x))$. Every nonlinear atom input
is a root strictly greater than $f(x)$. Equal peaks of two tags force equal
indices and equal argument values.
\end{lemma}
\begin{proof}
Write an atom norm as $A_0+bx+\sum a_jh_{\lambda_j}(A_j)$ as in \eqref{u:display:2-3}.
For an affine input, $b\ge1$ and
$x\le A_0+bx<f(x)$ once $x\ge\max(4,A_0+b+1)$.
For a nonlinear input the largest child peak $Y$ is at least $f(x)$;
for $x\ge\max(4,A_0+b,2C+1)$ the argument proving \eqref{u:display:2-4} applies
at every such $x$. It is a root greater than $f(x)$.

Two distinct inputs in $[x,f(x))$ cannot lie on the same $f$-orbit,
since the first positive iterate of the smaller is already at least
$f(x)$. Two nonlinear roots are handled by cancellation. In a mixed pair,
orbit cancellation either puts the larger nonlinear root in the range of
$f$, expresses the smaller affine input as an iterate of that larger
input, or identifies their disjoint bands; each is impossible. Thus equal
peaks force equal initial norms and equal iterate counts. Equation \eqref{u:display:1-3}
and controlled strictness then identify the two ordinals, exactly as in
Lemma~\ref{u:lem:roots}.
\end{proof}

\begin{theorem}[Effective finite zeros]\label{u:thm:zeros}
Given $E,F\in\mathcal P_1$, one computes $N\ge1$ such that either their
normal forms agree, or $E(x)\ne F(x)$ for every $x\ge N$. In the latter
case their equality set is the computable finite set obtained by exact
evaluation at $1,\ldots,N-1$.
\end{theorem}
\begin{proof}
Use the same finite set $\mathcal S$ as in Theorem~\ref{u:thm:identity}.
Take all bounds from Lemma~\ref{u:lem:oneinput}. For each pair in
$\mathcal S$, cancel its tag dictionaries. If a tag remains, include
$4$ and $2\sum|c_j|+|p|+|q|$ in the maximum; otherwise include
$|q|+1$ when its affine slope $p$ is nonzero. Height induction now
separates distinct argument forms and, by Lemma~\ref{u:lem:oneinput}, the
peaks of distinct tags. The sparse leading-peak estimate \eqref{u:display:2-5} separates
the full rows on this entire tail. These finitely many bounds are computed
before the induction and work simultaneously at every height.
\end{proof}

\begin{lemma}[Range cells and signed rows]\label{u:lem:cells}
The following assertions have effective tail bounds.
For a nonlinear positive expression $E$, let $Z_i$ be its displayed open
tag peaks, and $Y=\max_i Z_i=f(s)$. Then
\begin{equation}
              Y<E<f(s+1),\qquad E<h_\omega(Y).        \label{u:display:3-1}
\end{equation}
For a value $V\in\operatorname{range}(f)$ on that tail,
\begin{equation}
 V<E\ \Longleftrightarrow\ (\exists i\ Z_i\ge V),\qquad
 V>E\ \Longleftrightarrow\ (\forall i\ Z_i<V).       \label{u:display:3-2}
\end{equation}
After cancellation of identical tags in a signed row, its eventual
pointwise sign is the coefficient sign of its greatest remaining peak,
once those peaks have been ordered. If no tags remain, the row is affine.
\end{lemma}
\begin{proof}
Every open child uses the same input $x$, so its peak dominates the entire
fixed affine part on an effective tail. The proof of \eqref{u:display:2-4}, with the
outer index norm omitted, gives $Y<E<(2C+1)Y<f(s+1)$.
Increasing the threshold makes this fixed multiple smaller than
$h_\omega(Y)=f(Y+2)+1$. There is no range value strictly between $f(s)$
and $f(s+1)$, proving \eqref{u:display:3-2}, including impossibility of $V=E$.
Theorem~\ref{u:thm:zeros}'s height induction separates the remaining tag peaks,
and the estimate in \eqref{u:display:2-5} proves the signed-row assertion. It does not
identify the greatest peak; that task is supplied below.
\end{proof}

